\section{Chapter 4 supplementary exercises}
These exercises connect spanning sets, bases, coordinate changes, and rank.
When proving two sets equal, establish both inclusions. When proving a
universal claim false, give one concrete counterexample. For controllability,
form $[\,B\ AB\ \cdots\ A^{n-1}B\,]$ and check whether its columns span
all of the state space $\R^n$.

\exercise{4.S.1}{Twenty statements tying the chapter together}
\begin{enumerate}[label=(\alph*)]
\item \textbf{True.} The set is $\Span\{v_1,\ldots,v_p\}$, a subspace
of $V$ by the span theorem.
\item \textbf{True.} Adding another generator cannot remove combinations
already obtainable. A spanning subset remains spanning after a vector is added.
\item \textbf{False.} Start with the independent list $\{e_1\}$ in
$\R^2$ and add $e_1$ again; alternatively add $2e_1$. The enlarged list is dependent.
\item \textbf{False.} In $\R^2$, $\{e_1\}$ is independent but does not
span $\R^2$, so it is not a basis.
\item \textbf{True.} The Spanning Set Theorem removes any redundant
vectors from a finite spanning list, leaving an independent spanning subset.
\item \textbf{True.} A list of exactly $\dim V=p$ vectors that spans
$V$ is a basis by the Basis Theorem, hence is independent.
\item \textbf{False.} A plane not through zero is not a subspace.
For example the plane $z=1$ misses the zero vector.
\item \textbf{False.} In $A=[\,I_2\ I_2\,]$ the nonpivot columns
are $e_1,e_2$, which are independent. Each is a combination of pivot
columns, but that does not force a dependence among the nonpivot columns themselves.
\item \textbf{True.} Rows $(1,0),(1,0)$ have the relation $r_2-r_1=0$;
replacing row two by row two minus row one changes the list to
$(1,0),(0,0)$, which no longer has that same relation.
\item \textbf{False.} Elementary row operations replace $Ax=0$ by an
equivalent system, so they preserve the null space.
\item \textbf{False.} The matrix
$\left(\begin{smallmatrix}1&0\\2&0\end{smallmatrix}\right)$ has two
nonzero rows but rank one. Count nonzero rows only in an echelon form.
\item \textbf{False.} Nullity is the number of columns minus rank,
$n-k$, not generally $m-k$.
\item \textbf{True.} Row-equivalent matrices have the same row space,
so their row-space dimensions, and hence ranks, agree.
\item \textbf{False.} Nonzero rows need not be independent, as the matrix
in part (k) shows. Nonzero rows of an echelon form do give a row-space basis.
\item \textbf{True.} Row operations preserve row spaces, and both row
spaces equal that of the common reduced echelon form.
\item \textbf{True.} Choose a basis $h_1,\ldots,h_d$ for $H$, where
$d\le3$, and use these as columns of a $3\times3$ matrix, filling any
remaining columns with zeros. Its column space is $H$. If $H=\{0\}$,
use the zero matrix.
\item \textbf{False.} Rank $m$ means onto, not necessarily one-to-one.
For $A=[\,1\ 0\,]$, rank is $m=1$ but $(0,1)$ is a nonzero kernel vector.
\item \textbf{True.} Onto $\R^m$ means the column space has dimension
$m$, which is exactly rank $m$.
\item \textbf{True.} The reverse change-of-coordinates matrix is the
inverse, since changing from one basis to the other and back restores every vector.
\item \textbf{False.} Column $j$ of $P_{C\leftarrow B}$ is $[b_j]_C$,
not $[c_j]_B$. The latter columns occur in the inverse matrix.
\end{enumerate}

\exercise{4.S.2}{A parametrization has a redundant generator}
Collect the coefficients of $a,b,c$ to get generators
\[
 u=(1,2,-1,3),\quad v=(-2,5,-4,1),\quad w=(5,-8,7,1).
\]
But $w=u-2v$, so the third generator is redundant. The first two are
independent: if $\alpha u+\beta v=0$, the first two coordinates give
$\alpha-2\beta=0$, $2\alpha+5\beta=0$, hence $9\beta=0$ and both vanish.
\answer{A basis is $\{(1,2,-1,3),(-2,5,-4,1)\}$. There are three
parameters in the original expression, but only two independent directions.}

\exercise{4.S.3}{Describe a plane by a homogeneous equation}
Both $u_1=(-2,4,-6)$ and $u_2=(1,2,-5)$ satisfy
$x_1+2x_2+x_3=0$. They are not multiples, so their span is a two-dimensional
subspace contained in the plane defined by that equation.
To verify equality without relying only on dimensions, let $b$ satisfy
$b_1+2b_2+b_3=0$. Solve the first two equations of
$au_1+cu_2=b$: they give
\[
 a=(b_2-2b_1)/8,\qquad c=b_1+2a.
\]
With these choices the third coordinate is
$-6a-5c=-b_1-2b_2=b_3$.
\answer{$W=\{b\in\R^3:b_1+2b_2+b_3=0\}$.}

\exercise{4.S.4}{A scalar multiplier cannot depend on the input variable}
The equality $g(t)=tf(t)$ uses the function $t$, not one fixed real scalar.
Linear dependence would require constants $a,b$, not both zero, with
$a(3+t)+b(3t+t^2)=0$ for every $t$.
Comparing the $t^2$ coefficient gives $b=0$; comparing the constant gives
$3a=0$, so $a=0$ as well. Thus the pair is independent, and the proposed
discussion incorrectly treats a variable multiplier as a scalar.

\exercise{4.S.5}{Delete redundant polynomials from the given list}
The polynomials are $p_1=1+t$, $p_2=1-t$, $p_3=4$,
$p_4=t+t^2$, $p_5=1+2t+t^2$.
First $p_3=2p_1+2p_2$, so deleting $p_3$ does not alter the span.
Next $p_5=p_1+p_4$, so delete $p_5$.
The remaining $p_1,p_2,p_4$ are independent: in
$ap_1+bp_2+cp_4=0$, the quadratic coefficient gives $c=0$;
then the constant and linear coefficients give $a+b=0$ and $a-b=0$,
forcing $a=b=0$.
\answer{$\{1+t,1-t,t+t^2\}$ is a basis for $H$.}

\exercise{4.S.6}{A basis of a known two-dimensional span}
At least one of the four spanning polynomials is nonzero; choose it as $p$.
Since their span has dimension two, at least one other polynomial $q$ is
not a scalar multiple of $p$. Find such a $q$ by comparing coefficients.
Two nonproportional nonzero vectors are independent, so $\{p,q\}$ is an
independent pair in the two-dimensional space $H$. By the Basis Theorem
it is already a basis. There is no need to row reduce all four coefficient
vectors when the dimension is known in advance.

\exercise{4.S.7}{The homogeneous information that certifies every right-hand side}
For eighteen equations in twenty variables, every nonhomogeneous right-hand
side is attainable precisely when rank is eighteen. By rank--nullity this
is equivalent to nullity $20-18=2$. Thus it suffices, and is necessary, to
know that the homogeneous solution set is two-dimensional: exactly two
independent homogeneous solutions span it.

\exercise{4.S.8}{Equal finite dimensions force equality of nested spaces}
Let a basis of $H$ have $n$ vectors. Because $H\subseteq V$, these are
$n$ independent vectors in the $n$-dimensional space $V$.
The Basis Theorem says they span $V$. But their span is also $H$, so $H=V$.
The containment hypothesis matters: distinct subspaces can have equal
dimensions without either containing the other.

\exercise{4.S.9}{Rank and nullity of injective and surjective maps}
\step{(a) If $T$ is one-to-one.}
Its kernel is zero, so rank--nullity gives $\dim\operatorname{range}T=n$.
This also requires $n\le m$, since the range is contained in $\R^m$.
\step{(b) If $T$ maps onto $\R^m$.}
Its range has dimension $m$. Thus its kernel has dimension $n-m$,
which in particular requires $n\ge m$.

\exercise{4.S.10}{A maximal independent set must span}
Suppose the independent set $S$ failed to span $V$. Choose
$v\in V\setminus\Span S$. Then adjoining $v$ still leaves an independent
set: if $av+\sum_i c_i s_i=0$ is a finite relation and $a\ne0$, it would
express $v$ as a combination of elements of $S$, a contradiction.
Thus $a=0$, and independence of $S$ forces all $c_i=0$.
This contradicts maximality. Hence $S$ spans, and being independent as
well, it is a basis. The argument covers possibly infinite $S$ because
linear combinations always involve only finitely many terms.

\exercise{4.S.11}{A minimal finite spanning set is independent}
Suppose a finite spanning set were dependent. Then a nontrivial relation
would allow one vector with a nonzero coefficient to be solved for as a
combination of the others. Deleting that vector would preserve every
original linear combination, and thus preserve the span $V$.
That contradicts minimality. Therefore the set is independent; since it
already spans $V$, it is a basis.

\exercise{4.S.12}{Rank cannot increase beyond either factor's rank}
\step{(a) Compare column spaces.}
Every output $ABx=A(Bx)$ belongs to $\Col A$, so $\Col(AB)\subseteq\Col A$.
Taking dimensions gives $\rank(AB)\le\rank A$.
\step{(b) Apply that result to the transposes in reversed order.}
Using $(AB)^T=B^TA^T$ and transpose invariance of rank,
\[
 \rank(AB)=\rank(B^TA^T)\le\rank B^T=\rank B.
\]
Both inequalities follow. Repeated application proves the corresponding
bound for a product with more than two factors.

\exercise{4.S.13}{An invertible left factor preserves rank}
The product inequality gives $\rank(PA)\le\rank A$.
Since $A=P^{-1}(PA)$, applying it again gives
$\rank A\le\rank(PA)$. The two opposite inequalities imply equality.
Invertibility is what supplies the reverse factorization.

\exercise{4.S.14}{An invertible right factor preserves rank}
The product inequality first gives $\rank(AQ)\le\rank A$.
The identity $A=(AQ)Q^{-1}$ gives the reverse inequality
$\rank A\le\rank(AQ)$, hence equality. Equivalently one can apply
Exercise 4.S.13 to $(AQ)^T=Q^TA^T$.

\exercise{4.S.15}{A zero product limits two ranks}
If $AB=0$, every column of $B$ is in $\Nul A$.
Thus $\Col B\subseteq\Nul A$, and dimensions give
\[
 \rank B\le\dim\Nul A=n-\rank A.
\]
Rearrangement yields $\rank A+\rank B\le n$.
The shared $n$ is the number of inputs of $A$ and outputs of $B$,
which makes the containment meaningful.

\exercise{4.S.16}{Rank is subadditive}
Take rank factorizations $A=C_AR_A$, $B=C_BR_B$, with ranks $r,s$.
Then block multiplication gives
\[
 A+B=[\,C_A\ C_B\,]\begin{pmatrix}R_A\\R_B\end{pmatrix}.
\]
The left block matrix has $r+s$ columns and therefore rank at most $r+s$.
By the product inequality, $\rank(A+B)\le r+s=\rank A+\rank B$.
If either matrix has rank zero, it is zero and the inequality is immediate,
so the proof also covers those degenerate cases.

\exercise{4.S.17}{Find an invertible square submatrix of full rank}
\step{(a) Select independent columns.}
Choose the $r$ pivot columns of $A$ to obtain an $m\times r$ matrix $A_1$.
They are independent, so $\rank A_1=r$.
\step{(b) Select independent rows of that submatrix.}
Its row rank also equals $r$, so choose $r$ independent rows of $A_1$.
The resulting $r\times r$ matrix $A_2$ has $r$ independent rows, hence
rank $r$ and is invertible by the Invertible Matrix Theorem.
This proves the requested existence; the case $r=0$ has the usual empty
square-matrix convention or may be excluded as trivial.

\exercise{4.S.18}{Expand four controlled time steps}
\step{(a) Compute each state before collecting the blocks.}
From $x_0=0$ and $x_{k+1}=Ax_k+Bu_k$,
\begin{align*}
 x_1&=Bu_0,\\
 x_2&=ABu_0+Bu_1,\\
 x_3&=A^2Bu_0+ABu_1+Bu_2,\\
 x_4&=A^3Bu_0+A^2Bu_1+ABu_2+Bu_3\\
 &=\underbrace{[\,B\ AB\ A^2B\ A^3B\,]}_{M}
    \begin{pmatrix}u_3\\u_2\\u_1\\u_0\end{pmatrix}.
\end{align*}
The order of the input blocks is reversed to match the stated order in $M$.
Each $u_j$ has two entries, so the stacked input has eight entries,
matching the $4\times8$ matrix $M$.
\step{(b) Use full row rank to realize an arbitrary target.}
Controllability means rank $M=4$, so $M$ maps $\R^8$ onto $\R^4$.
Given $v\in\R^4$, solve $Mz=v$, then split its eight entries into four
two-entry blocks $(u_3,u_2,u_1,u_0)$. The computation in part (a) shows
that this control sequence gives $x_4=v$.

\exercise{4.S.19}{Check the rank of the controllability matrix}
Here the state dimension is $3$ and the input has one coordinate, so form $M=[\,B\ AB\ A^2B\,]$. Exact arithmetic gives
\[ M=\begin{pmatrix}0&1&0\\1&- \frac{9}{10}&\frac{81}{100}\\1&\frac{1}{2}&\frac{1}{4}\end{pmatrix},\qquad\det M=\frac{14}{25}.\]
Since this determinant is nonzero, $M$ has rank $3$ and its columns span the state space. Therefore the pair $(A,B)$ is controllable.

\exercise{4.S.20}{Check the rank of the controllability matrix}
Here the state dimension is $3$ and the input has one coordinate, so form $M=[\,B\ AB\ A^2B\,]$. Exact arithmetic gives
\[ M=\begin{pmatrix}1&\frac{1}{2}&\frac{19}{100}\\1&\frac{7}{10}&\frac{9}{20}\\0&0&0\end{pmatrix},\qquad\det M=0.\]
The determinant is zero, so $M$ does not have full rank and $(A,B)$ is not controllable. Its actual rank is $2$. Its last row is zero, while the leading $2\times2$ minor has determinant $1/5\ne0$, verifying rank two.

\exercise{4.S.21}{Check the rank of the controllability matrix}
Here the state dimension is $4$ and the input has one coordinate, so form $M=[\,B\ AB\ A^2B\ A^3B\,]$. Exact arithmetic gives
\[ M=\begin{pmatrix}1&0&0&-1\\0&0&-1&\frac{8}{5}\\0&-1&\frac{8}{5}&- \frac{24}{25}\\-1&\frac{8}{5}&- \frac{24}{25}&- \frac{3}{125}\end{pmatrix},\qquad\det M=0.\]
The determinant is zero, so $M$ does not have full rank and $(A,B)$ is not controllable. Its actual rank is $3$. For example its leading $3\times3$ minor has determinant $-1$, showing rank at least three; the zero full determinant then gives exactly three.

\exercise{4.S.22}{Check the rank of the controllability matrix}
Here the state dimension is $4$ and the input has one coordinate, so form $M=[\,B\ AB\ A^2B\ A^3B\,]$. Exact arithmetic gives
\[ M=\begin{pmatrix}1&0&0&-1\\0&0&-1&\frac{1}{2}\\0&-1&\frac{1}{2}&\frac{229}{20}\\-1&\frac{1}{2}&\frac{229}{20}&- \frac{411}{40}\end{pmatrix},\qquad\det M=- \frac{3}{10}.\]
Since this determinant is nonzero, $M$ has rank $4$ and its columns span the state space. Therefore the pair $(A,B)$ is controllable.
